\chapter{PENDAHULUAN}

\section{Latar Belakang}

% TODO: Uraikan latar belakang secara berurutan:
% (1) persoalan nyata yang akan diselesaikan, disertai alasan, urgensi, dan
%     fenomena nyata yang memotivasi penelitian;
% (2) apa yang telah dilakukan peneliti sebelumnya untuk mengatasi persoalan
%     tersebut, dengan mengacu pada karya ilmiah (laporan TA, publikasi
%     nasional/internasional);
% (3) solusi yang ditawarkan melalui metode, pendekatan, teknik, atau strategi
%     tertentu;
% (4) target yang akan dicapai dalam penelitian.
% Hubungkan pula roadmap riset/teknologi penelitian dengan roadmap departemen,
% fakultas, atau universitas.

Lorem ipsum dolor sit amet, consectetur adipiscing elit. Sed do eiusmod tempor incididunt ut labore et dolore magna aliqua. Ut enim ad minim veniam, quis nostrud exercitation ullamco laboris nisi ut aliquip ex ea commodo consequat. Duis aute irure dolor in reprehenderit in voluptate velit esse cillum dolore eu fugiat nulla pariatur \citep{referensi1}.

Sed ut perspiciatis unde omnis iste natus error sit voluptatem accusantium doloremque laudantium, totam rem aperiam, eaque ipsa quae ab illo inventore veritatis et quasi architecto beatae vitae dicta sunt explicabo \citep{referensi2}. Nemo enim ipsam voluptatem quia voluptas sit aspernatur aut odit aut fugit, sed quia consequuntur magni dolores eos qui ratione voluptatem sequi nesciunt.

\section{Perumusan Masalah}

% TODO: Tuliskan rumusan masalah dalam bentuk deklaratif atau pernyataan yang
% tegas dan jelas (bukan pertanyaan penelitian dan bukan akar permasalahan).
% Rumusan masalah merupakan kesenjangan antara keadaan yang ada dengan keadaan
% yang akan dicapai, umumnya berkaitan dengan kinerja sistem, aplikasi, atau
% fungsionalitas sistem.

Berdasarkan latar belakang di atas, permasalahan penelitian ini dirumuskan sebagai berikut:

\begin{enumerate}
  \item Lorem ipsum dolor sit amet, consectetur adipiscing elit.
  \item Sed do eiusmod tempor incididunt ut labore et dolore magna aliqua.
  \item Ut enim ad minim veniam, quis nostrud exercitation ullamco laboris.
\end{enumerate}

\section{Batasan Penelitian}

% TODO: Tuliskan batasan penelitian untuk menjaga fokus dan menghindari
% perluasan yang tidak perlu. Batasan dapat mencakup ruang lingkup penelitian,
% asumsi, data atau dataset yang digunakan, serta lokasi penelitian atau
% lokasi pengambilan data.

Agar penelitian tetap terfokus dan terukur, batasan penelitian yang ditetapkan adalah:

\begin{enumerate}
  \item Lorem ipsum dolor sit amet, consectetur adipiscing elit.
  \item Sed do eiusmod tempor incididunt ut labore et dolore magna aliqua.
  \item Ut enim ad minim veniam, quis nostrud exercitation ullamco laboris.
\end{enumerate}

\section{Tujuan Penelitian}

% TODO: Tuliskan tujuan penelitian yang menggambarkan apa yang akan dikerjakan
% dan hasil yang diharapkan dalam rangka menyelesaikan permasalahan pada
% rumusan masalah. Tujuan penelitian menjadi target selesainya penelitian.

Tujuan dari penelitian ini adalah:

\begin{enumerate}
  \item Lorem ipsum dolor sit amet, consectetur adipiscing elit.
  \item Sed do eiusmod tempor incididunt ut labore et dolore magna aliqua.
  \item Ut enim ad minim veniam, quis nostrud exercitation ullamco laboris.
\end{enumerate}

\section{Manfaat Penelitian}

% TODO: Tuliskan manfaat penelitian bagi dunia nyata dan/atau lingkungan ilmiah.

Penelitian ini diharapkan memberikan manfaat berikut:

\begin{enumerate}
  \item \textbf{Teoritis:} Lorem ipsum dolor sit amet, consectetur adipiscing elit.
  \item \textbf{Praktis:} Sed do eiusmod tempor incididunt ut labore et dolore magna aliqua.
\end{enumerate}
